%-------------------------------------------------------------------------------
% 7. 2025 제2회 MixUP AI Datathon — 진동/음향 데이터 회귀 예측 (Track 1 출제위원)
%-------------------------------------------------------------------------------
\cvsection{8. 2025 제2회 MixUP AI Datathon : 진동/음향 데이터 회귀 예측}

\projsummary{산업 현장 비정형 데이터를 활용한 '누수 수준 예측' 회귀 과제 설계 및 대회 운영 --- Track 1 총괄 운영·출제위원}
\projfigL[5.4cm]{p7.jpg}

\projpart{\faBullhorn\hspace{1.3mm}배경 및 목표}
\begin{cvitems}
  \item {\textbf{배경}\enskip 4개 연합 동아리(ToBig's·BITAmin·Prometheus·KUBIG)가 공동 주최한 제2회 연합 데이터톤.}
  \item {\textbf{목표}\enskip NIA 제공 데이터로 누수 수준을 예측하는 회귀(Regression) 과제를 기획하고, 참가자가 공정하게 경쟁할 Kaggle 기반 평가 환경을 구축.}
  \item {\textbf{운영 기간}\enskip 2025.11.22 -- 2025.11.23 (준비 기간 약 3개월)}
  \item {\textbf{기술 스택}\enskip \pill{Python} \pill{Pandas} \pill{Scikit-learn} \pill{Kaggle API} \pill{Metric Optimization}}
\end{cvitems}

\projpart{\faChartBar\hspace{1.3mm}데이터 소개}
\begin{cvitems}
  \item {\textbf{데이터셋}\enskip NIA(한국지능정보사회진흥원) 제공 상수관로 고해상도 진동/음향 센서 로그.}
  \item {\textbf{누수 유형(leaktype)}\enskip in(옥내)·out(옥외)·other(기타)·noise(노이즈)·normal(정상) 5종.}
  \item {\textbf{누수 수준(llevel)}\enskip 누수 정도를 나타내는 연속형 타겟 변수.}
  \item {\textbf{주파수 피처}\enskip 음향·진동 신호를 푸리에 변환한 513차원(0HZ--5120HZ) 대역별 강도.}
\end{cvitems}

\projpart{\faMagnifyingGlass\hspace{1.3mm}설계 과정 및 수행 내용}

\stephead{1}{변별력 확보를 위한 회귀 과제 설계}
\begin{substeps}
  \item 단순 누수 유형 분류는 특정 관로 고유 패턴을 암기하는 데이터 누수 리스크에 취약.
  \item 과제를 분류가 아닌 누수 정도(llevel) 정밀 예측 회귀로 설계해 변별력 확보.
  \item 유형별 타겟 분포 왜곡을 막기 위해 50\% 층화 추출(Stratified Sampling)로 학습셋 구축.
\end{substeps}

\stephead{2}{데이터 무결성 및 부정행위(Leakage) 차단}
\begin{substeps}
  \item 동일 장소 신호 혼입 시 모델이 물리 특징이 아닌 장소 패턴을 암기하고, 식별자(ID)가 정답 힌트가 될 수 있음.
  \item 정답 유추 가능한 장소 식별자(sid·site)와 누수율(lrate)을 제거.
  \item 계절성 피처용 날짜(ldate)는 보존하고 중복 정보인 MAX 관련 컬럼은 사전 정제해 배포.
  \item 관로 ID별 데이터 분리 로직으로 미학습 장소의 누수를 예측하게 해 평가 신뢰도 확보.
\end{substeps}

\stephead{3}{평가 지표 최적화 및 강건한 학습 환경}
\begin{substeps}
  \item 핵심 변수(llevel·ldate) 제외 피처에 5\% 마스킹을 적용해 주파수 전체 상관관계를 학습하도록 유도.
  \item 전체 평가 데이터를 Public(30\%) / Private(70\%)로 나누고 최종 순위는 Private로만 결정.
  \item 초기 MSLE의 낮은 해상도를 확인하고, 모델 설명력을 직관적으로 반영하는 $R^2$ Score(0.92)로 지표를 도입해 변별력 확보.
\end{substeps}

\achievement{대회 이슈 \textbf{0건 달성} --- 3개월간 과제 설계·평가 환경을 총괄 운영.}
